\section{Introduction}
\label{sec:introduction}

When a careful writer composes a paragraph, she does not simply transcribe an unbroken stream of thought.
She writes a sentence, pauses to consider whether it is accurate, checks that it follows logically from the previous one, and only then moves on.
Large language models, by contrast, produce each token conditioned only on the tokens already generated, with no explicit mechanism for reconsidering what comes next at the level of a complete sentence.
Recent work on \emph{pause tokens}~\citep{goyal2024think} and \emph{test-time compute scaling}~\citep{snell2024scaling,muennighoff2025s1} has shown that giving models additional computation at inference time can improve task performance.
However, these approaches either require training modifications or operate at the response level, and none study how the \emph{character} of the generated text changes when models are forced to deliberate sentence by sentence.

We ask a simple question: {\bf what happens to an LLM's writing when it must think before every sentence?}
To answer this, we introduce \iscot (\iscotshort), a prompting strategy that instructs the model to produce a \texttt{<think>} block of explicit reasoning before each sentence of its visible response.
Unlike conventional chain-of-thought~\citep{wei2022chain}, which places all reasoning before the final answer, \iscotshort interleaves reasoning with output at the sentence level, creating a pattern in which every claim is preceded by deliberation about its accuracy, phrasing, and logical connection to the surrounding text.

We evaluate \iscotshort on \gptfour across three tasks: 30 open-ended generation prompts, 100 \truthfulqa questions, and 100 \gsmk math problems, comparing it against standard generation and standard CoT.
Our automated text analysis reveals that \iscotshort produces a distinctive ``careful writer'' style: sentences are longer and more information-dense (20.9 vs.\ 15.4 words per sentence, $d = 1.07$), lexical diversity is significantly higher (type-token ratio 0.700 vs.\ 0.640, $d = 0.84$), and hedging and self-correction language appears more frequently.
On \truthfulqa, \iscotshort achieves the highest truthfulness rate (98\% vs.\ 95\% for standard generation), with qualitative evidence that deliberation helps avoid false presuppositions.
On \gsmk, all three conditions perform equivalently (90--91\%).

The results also reveal an unexpected trade-off.
Because \texttt{<think>} blocks consume part of the model's token budget, \iscotshort responses are substantially shorter (130.8 vs.\ 200.6 visible words).
LLM-as-judge evaluations overwhelmingly prefer the longer standard and CoT responses, highlighting a well-known length bias in automated evaluation~\citep{zheng2023judging} that obscures the per-sentence quality improvements \iscotshort achieves.

We make the following contributions:
\begin{itemize}[leftmargin=*,itemsep=0pt,topsep=0pt]
    \item We propose \iscot, a prompting strategy that requires explicit reasoning before every sentence, and show that it produces measurably different text without any model training.
    \item We conduct the first systematic study of how sentence-level deliberation affects text properties---hedging, lexical diversity, sentence structure, and self-correction---across open-ended, factual, and mathematical tasks.
    \item We identify a token-budget trade-off between per-sentence quality and response breadth, and demonstrate that standard LLM-as-judge evaluation is poorly calibrated to detect the quality shift that \iscotshort produces.
\end{itemize}
